\chapter{Arquitetura}
\label{cap:arquitetura}

\section{Vis\~ao em camadas}

O kernel \'e organizado em seis camadas conceituais. As camadas inferiores
dependem apenas de hardware e de \code{core}; as superiores usam as
inferiores. N\~ao h\'a \code{std}, alocador de heap, \emph{runtime} de Rust
nem bibliotecas de terceiros.

\begin{figure}[H]
\centering
\begin{tikzpicture}[
  layer/.style={draw=rinuxblue, fill=rinuxlight, rounded corners=2pt,
                minimum width=11.4cm, minimum height=1.0cm,
                align=center, font=\small},
  hw/.style={draw=black!60, fill=black!8, rounded corners=2pt,
             minimum width=11.4cm, minimum height=1.0cm,
             align=center, font=\small},
  node distance=0.35cm
]
\node[layer] (shell) {\textbf{Interface} --- shell: linha de comando, \code{help}, \code{echo}, \code{info}, \code{demand}, \code{usertest}};
\node[layer, below=of shell] (services) {\textbf{Servi\c{c}os} --- console VGA, entrada PS/2, serial, macros de impress\~ao};
\node[layer, below=of services] (mem) {\textbf{Mem\'oria e interrup\c{c}\~oes} --- GDT/TSS, IDT/PIC/PIT, frames e \emph{demand paging}};
\node[layer, below=of mem] (core) {\textbf{N\'ucleo} --- \code{kernel\_main}, panic handler, spinlock, portas de I/O};
\node[layer, below=of core] (boot) {\textbf{Boot} --- stub 32/64 bits, GDT inicial, PML4/PDPT/PD, Multiboot2};
\node[hw, below=of boot] (hw) {\textbf{Hardware} --- CPU x86\_64, VGA, 8042, UART 16550, PIC/PIT, BIOS/GRUB};
\draw[-{Latex[length=2mm]}, thick, rinuxblue] (shell) -- (services);
\draw[-{Latex[length=2mm]}, thick, rinuxblue] (services) -- (mem);
\draw[-{Latex[length=2mm]}, thick, rinuxblue] (mem) -- (core);
\draw[-{Latex[length=2mm]}, thick, rinuxblue] (core) -- (boot);
\draw[-{Latex[length=2mm]}, thick, black!60] (boot) -- (hw);
\end{tikzpicture}
\caption{Camadas do Rust VGA Kernel.}
\label{fig:camadas}
\end{figure}

\section{Modelo de execu\c{c}\~ao}

O kernel executa em \textbf{ring 0}, em uma \'unica CPU, com
\textbf{interrup\c{c}\~oes desligadas}. Todo o processamento \'e feito por
\emph{polling}: o shell consulta o teclado e a serial em um la\c{c}o ativo,
usando \code{core::hint::spin\_loop()} quando n\~ao h\'a entrada.

\begin{table}[H]
\centering
\caption{Caracter\'isticas do modelo de execu\c{c}\~ao.}
\label{tab:modelo}
\begin{tabularx}{\linewidth}{@{}lX@{}}
\toprule
\textbf{Aspecto} & \textbf{Implementa\c{c}\~ao} \\
\midrule
Privil\'egio & Ring 0 para o kernel; ring 3 apenas nos demos de
\code{user.rs}, com GDT/TSS. \\
CPUs & Uma CPU; SMP n\~ao \'e inicializado. \\
Interrup\c{c}\~oes & IDT com 256 portas; PIC remapeado e timer IRQ0 a
~100 Hz. IRQ1 (teclado) permanece mascarada. \\
Entrada & Polling de \code{0x60}/\code{0x64} e de COM1; timer por IRQ. \\
Concorr\^encia & Spinlocks seguros para IRQ no console, ramfs, heap, VFS e
tabela de tarefas. O escalonador \emph{round-robin} preempta no quantum do
timer (5 ticks), troca em yield expl\'icito, bloqueia tarefas em \code{sleep}
at\'e um tick de prazo e as suspende em filas de espera at\'e um
\code{wake\_up}. \\
Mem\'oria & Frames f\'isicos de 4 KiB, tabelas sob demanda e page faults, um
heap \emph{first-fit} de 1 MiB e \code{alloc}/\code{Vec}/\code{String}. \\
Pilhas & 8 KiB no stub (baixa), 32 KiB no kernel (alta), 16 KiB para
interrup\c{c}\~oes de usu\'ario e 8 KiB de IST para double fault. \\
Tempo de vida & O shell nunca retorna; \code{halt} e \code{reboot} encerram o fluxo. \\
\bottomrule
\end{tabularx}
\end{table}

\section{Depend\^encias entre m\'odulos}

\begin{table}[H]
\centering
\caption{Depend\^encias principais entre m\'odulos.}
\label{tab:dependencias}
\begin{tabularx}{\linewidth}{@{}lX@{}}
\toprule
\textbf{M\'odulo} & \textbf{Usa / \'e usado por} \\
\midrule
\code{main.rs} & Declara todos os m\'odulos; chama \code{vga}, \code{serial},
\code{tty}, \code{paging}, \code{heap}, \code{ramfs}, \code{vfs}, \code{gdt},
\code{keyboard}, \code{idt}, \code{task}, \code{arch} e \code{shell};
implementa as macros de impress\~ao. \\
\code{boot.rs} & Assembly que chama \code{kernel\_main}; n\~ao \'e chamado por
Rust. \\
\code{arch.rs} & Portas de I/O, \code{irq\_save}/\code{irq\_restore},
\code{sti}, \code{halt} e \code{reboot}; usado por v\'arios m\'odulos. \\
\code{gdt.rs} & GDT/TSS e pilhas; usa apenas \code{core::arch}. \\
\code{idt.rs} & Usa \code{arch}, \code{paging}, \code{syscall} e
\code{user} para tratar exce\c{c}\~oes, IRQs e a \emph{syscall}. \\
\code{paging.rs} & Frames, tabelas e \emph{page faults}; usa apenas
\code{core::arch}. \\
\code{heap.rs} & Heap \emph{first-fit} e \code{\#[global\_allocator]}; usa
\code{paging} e \code{sync}. \\
\code{ramfs.rs} & \'Arvore em mem\'oria de 48 n\'os; usa \code{sync}. \\
\code{proc.rs} & \code{/proc} sint\'etico; renderiza a partir de \code{heap},
\code{idt}, \code{paging} e \code{task}. \\
\code{vfs.rs} & Tabela de descritores sobre ramfs, \code{/proc} e consoles;
usa \code{proc}, \code{ramfs}, \code{tty} e \code{sync}. \\
\code{task.rs} & Tabela de tarefas, pilhas, troca preemptiva por timer,
\code{sleep} bloqueante e filas de espera por evento; usa \code{sync} e
\code{arch}. \\
\code{user.rs} & Monta os programas, chama \code{paging} e \code{tty} e
retorna ao kernel por \code{user\_resume}. \\
\code{syscall.rs} & Implementa as onze chamadas; usa \code{paging},
\code{task}, \code{tty}, \code{user} e \code{vfs}. \\
\code{shell.rs} & Usa \code{tty}, \code{idt}, \code{paging}, \code{heap},
\code{ramfs}, \code{vfs}, \code{proc}, \code{task}, \code{user} e
\code{arch}. \\
\code{tty.rs} & Mant\'em os oito consoles; usa \code{vga}, \code{keyboard},
\code{serial} e \code{sync}. \\
\code{vga.rs} & Renderizador de baixo n\'ivel; usa \code{arch}. \\
\code{keyboard.rs} & Usa \code{arch} para acessar o 8042. \\
\code{serial.rs} & Usa \code{arch} para acessar a UART. \\
\code{sync.rs} & Spinlocks seguros para IRQ no console, ramfs, heap, VFS e
tarefas; usa \code{core} e \code{arch}. \\
\bottomrule
\end{tabularx}
\end{table}

A tabela mostra apenas as depend\^encias relevantes. \code{boot.rs} \'e um
caso especial: ele cont\'em o c\'odigo de montagem que chama \code{kernel\_main}
e n\~ao \'e chamado por Rust. \code{main.rs} declara todos os m\'odulos e
exporta as macros \code{print!} e \code{println!}.

\section{Cadeia de boot}

\begin{figure}[H]
\centering
\begin{tikzpicture}[
  stage/.style={draw=rinuxblue, fill=rinuxlight, rounded corners=3pt,
                minimum width=3.0cm, minimum height=1.25cm,
                align=center, font=\scriptsize},
  arrow/.style={-{Latex[length=2mm]}, thick, rinuxblue},
  node distance=0.55cm
]
\node[stage] (bios) {BIOS\\SeaBIOS};
\node[stage, right=of bios] (eltorito) {El Torito\\CD};
\node[stage, right=of eltorito] (grub) {GRUB\\i386-pc};
\node[stage, right=of grub] (mb) {Multiboot2\\kernel.elf};
\node[stage, below=1.0cm of mb] (stub) {\_start\\32 bits @ 1 MiB};
\node[stage, left=of stub] (paging) {PML4/PDPT/PD\\GDT + long mode};
\node[stage, left=of paging] (main) {kernel\_main\\metade alta};
\node[stage, left=of main] (shell) {vga / tty / kbd\\shell};
\draw[arrow] (bios) -- (eltorito);
\draw[arrow] (eltorito) -- (grub);
\draw[arrow] (grub) -- (mb);
\draw[arrow] (mb) -- (stub);
\draw[arrow] (stub) -- (paging);
\draw[arrow] (paging) -- (main);
\draw[arrow] (main) -- (shell);
\end{tikzpicture}
\caption{Cadeia de boot completa.}
\label{fig:bootchain}
\end{figure}

\section{Decis\~oes de projeto}

\begin{table}[H]
\centering
\caption{Decis\~oes arquiteturais e justificativas.}
\label{tab:decisoes}
\begin{tabularx}{\linewidth}{@{}p{3.3cm}X@{}}
\toprule
\textbf{Decis\~ao} & \textbf{Justificativa} \\
\midrule
\code{no\_std} & Elimina depend\^encia de sistema operacional, alocador e
\emph{runtime}; o c\'odigo usa apenas \code{core}. \\
Sem depend\^encias & Reduz superf\'icie de ataque e mant\'em o build
reprodut\'ivel; \code{Cargo.lock} n\~ao tem pacotes externos. \\
Alvo \code{x86\_64-unknown-none} & Alvo oficial bare metal, com modelo de
c\'odigo \emph{kernel} e sem biblioteca padr\~ao. \\
Kernel na metade alta & Conven\c{c}\~ao do Linux e do alvo; evita depender de
endere\c{c}os baixos e deixa o stub separado. \\
Stub em assembly & O controle inicial vem em 32 bits com pilha indefinida;
a transi\c{c}\~ao para 64 bits exige c\'odigo posicional e controle exato. \\
Tabelas de p\'agina est\'aticas & O stub n\~ao tem alocador; 3 p\'aginas de
4 KiB s\~ao suficientes para mapear 1 GiB com \emph{huge pages} de 2 MiB.
Depois do boot, \code{paging.rs} aloca tabelas de 4 KiB para as regi\~oes sob
demanda. \\
Polling no teclado & Mant\'em IRQ1 mascarada e evita um buffer de scancodes;
o timer IRQ0 j\'a demonstra o caminho de interrup\c{c}\~ao. \\
Serial al\'em do VGA & Permite depura\c{c}\~ao e testes automatizados sem
depender de captura gr\'afica. \\
\emph{Panic abort} & N\~ao h\'a \emph{unwinding} nem runtime de p\^anico;
\code{panic=abort} reduz c\'odigo gerado. \\
LTO e \code{codegen-units=1} & Gera um bin\'ario pequeno e previs\'ivel; o
linker descarta se\c{c}\~oes n\~ao usadas com \code{--gc-sections}. \\
\bottomrule
\end{tabularx}
\end{table}

\section{Tratamento de erros e falhas}

O kernel n\~ao possui um subsistema de erros. As estrat\'egias s\~ao:

\begin{itemize}
  \item \textbf{Timeout em esperas de hardware}: as fun\c{c}\~oes do 8042
        tentam por um n\'umero limitado de itera\c{c}\~oes e desistem, evitando
        \emph{hang} quando o controlador n\~ao responde.
  \item \textbf{Verifica\c{c}\~ao de presen\c{c}a da UART}: se o registrador de
        status devolver \code{0xFF}, a serial \'e considerada ausente.
  \item \textbf{Entrada inv\'alida}: bytes fora da faixa imprim\'ivel s\~ao
        ignorados no shell; UTF-8 inv\'alido em \code{echo} vira uma
        mensagem fixa.
  \item \textbf{Panic}: o handler imprime \code{[PANIC]} no console e entra em
        \code{halt\_loop()}. Exce\c{c}\~oes de CPU passam pela IDT; um
        \code{\#PF} sob demanda pode ser resolvido, mas uma falha fora das
        regi\~oes mapeadas imprime diagn\'ostico e para a CPU. Um
        \code{\#DF} usa a pilha IST e tamb\'em para o sistema.
  \item \textbf{Sem recupera\c{c}\~ao}: n\~ao h\'a \emph{kernel oops},
        reinicializa\c{c}\~ao controlada ou despejo de pilha.
\end{itemize}

\section{Concorr\^encia e sincroniza\c{c}\~ao}

Apesar de o sistema ser monoprocessado e rodar com interrup\c{c}\~oes
desligadas, o console \'e protegido por um \code{SpinLock}. Isso permite que
chamadas de impress\~ao vindas de contextos diferentes no futuro n\~ao
intercalem bytes no buffer VGA. O lock usa \code{AtomicBool::swap} com
ordena\c{c}\~ao \code{Acquire} na aquisi\c{c}\~ao e \code{Release} na
libera\c{c}\~ao. A serial n\~ao \'e protegida pelo lock; como o console
escreve VGA e serial em sequ\^encia, uma futura implementa\c{c}\~ao
multithread precisaria de um lock adicional para a UART.

\section{Ciclo de vida da inicializa\c{c}\~ao}

\begin{table}[H]
\centering
\caption{Ordem de inicializa\c{c}\~ao em \code{kernel\_main}.}
\label{tab:init}
\begin{tabular}{@{}clp{7.5cm}@{}}
\toprule
\textbf{\#} & \textbf{Chamada} & \textbf{Efeito} \\
\midrule
1 & \code{vga::init()} & Limpa a tela, zera o cursor, define forma do cursor. \\
2 & \code{serial::init()} & Configura COM1 em 38400 8N1 e detecta presen\c{c}a. \\
3 & \code{keyboard::init()} & Configura o 8042, tradu\c{c}\~ao e varredura. \\
4 & \code{println!} & Banner e mensagem de boas-vindas (VGA + serial). \\
5 & valida\c{c}\~ao & Confere \code{magic == 0x36d76289}. \\
6 & \code{shell::run()} & La\c{c}o eterno de prompt e leitura de comandos. \\
\bottomrule
\end{tabular}
\end{table}
